\chapter{Réalisation}\label{chap5}
\pagestyle{fancy}
\lhead{Chapitre 5 : Réalisation}

% =====================================================================
\section*{Introduction}
\addcontentsline{toc}{section}{Introduction}

Ce chapitre présente la réalisation de la plateforme VigiloChain : les outils
utilisés, l'API backend, le tableau de bord, la vérification par les tests et
le déploiement. Il se termine par une matrice de traçabilité qui relie chaque
besoin fonctionnel du chapitre~\ref{chap2} à la page du tableau de bord et aux
routes de l'API qui le réalisent.

% =====================================================================
\section{Environnement de travail}

Le tableau~\ref{tab:environnement} présente les principaux outils utilisés.

\begin{table}[H]
\centering
\small
\renewcommand{\arraystretch}{1.3}
\setlength{\arrayrulewidth}{0.8pt}
\arrayrulecolor{thead}
\caption{Environnement logiciel du projet VigiloChain.}
\label{tab:environnement}
\begin{tabularx}{\textwidth}{@{}>{\bfseries}p{3.6cm} X >{\centering\arraybackslash}p{3.2cm}@{}}
\rowcolor{thead}
\textcolor{white}{\textbf{Catégorie}} & \textcolor{white}{\textbf{Logiciel}} &
\textcolor{white}{\textbf{Version}} \\
\hline
Développement & Visual Studio Code, Jupyter Notebook & -- \\
\rowcolor{trowa}
Backend & Python, FastAPI, Pydantic & 3.12, 0.110+, 2.6+ \\
Données & PostgreSQL, SQLite, SQLModel, SQLAlchemy & 16, --, 0.0.24, 2.0+ \\
\rowcolor{trowa}
Apprentissage & scikit-learn, XGBoost, SHAP & 1.4+, 2.0+, 0.45+ \\
Frontend & React, Vite, TypeScript, React Router & 19, 8, 6, 7 \\
\rowcolor{trowa}
Interface & Tailwind CSS, Recharts, Radix UI, Lucide & 4, 3, --, -- \\
Sécurité & python-jose (JWT), bcrypt & -- \\
\rowcolor{trowa}
Déploiement & Docker, Docker Compose & -- \\
Qualité & pytest & 8+ \\
\hline
\end{tabularx}
\end{table}

% =====================================================================
\section{Réalisation du backend}

\subsection{Organisation du code}

Le backend est une application FastAPI organisée en quatre niveaux, de
l'extérieur vers l'intérieur : \texttt{api/v1} (routes REST), \texttt{domains}
(logique métier par domaine : \textit{auth}, \textit{transmission},
\textit{parametres}), \texttt{models} et \texttt{db} (entités SQLModel, session
asynchrone, données initiales). Une route ne contient que la validation des
entrées, le contrôle d'accès et l'appel au service du domaine concerné.

Au démarrage, l'application initialise la base, crée le compte administrateur
à partir des variables d'environnement et sert le tableau de bord compilé sur
le même port que l'API. Si PostgreSQL n'est pas joignable, elle bascule sur
SQLite.

\subsection{Authentification et contrôle d'accès}

La connexion (\texttt{POST /auth/login}) vérifie le mot de passe haché avec
bcrypt et renvoie un jeton JWT qui porte l'identité et le rôle. Chaque route
protégée déclare une dépendance FastAPI : \texttt{get\_current\_user} pour tout
utilisateur connecté, \texttt{require\_any\_role("operator")} pour l'opérateur
et l'administrateur, \texttt{require\_any\_role()} sans argument pour
l'administrateur seul. Le refus renvoie le code HTTP~403. Les actions
d'administration (comptes, mots de passe, paramètres) écrivent une ligne dans
la table \texttt{audit\_logs}.

\subsection{Routes de l'API}

Toutes les routes sont préfixées par \texttt{/api/v1}. Le
tableau~\ref{tab:api} liste celles que le tableau de bord appelle,
regroupées par module. Une documentation interactive (Swagger UI) est
générée automatiquement par FastAPI.

\begin{table}[H]
\centering
\footnotesize
\renewcommand{\arraystretch}{1.35}
\setlength{\arrayrulewidth}{0.8pt}
\arrayrulecolor{thead}
\caption{Routes de l'API utilisées par le tableau de bord.}
\label{tab:api}
\begin{tabularx}{\textwidth}{@{}>{\raggedright\arraybackslash}p{1.5cm} >{\raggedright\arraybackslash}p{5.5cm} >{\centering\arraybackslash}p{1.9cm} >{\raggedright\arraybackslash}X@{}}
\rowcolor{thead}
\textcolor{white}{\textbf{Méthode}} & \textcolor{white}{\textbf{Chemin}} &
\textcolor{white}{\textbf{Accès}} & \textcolor{white}{\textbf{Rôle de la route}} \\
\hline
\multicolumn{4}{@{}l}{\textit{Module auth}} \\
POST & \texttt{/auth/register} & Public & Inscription (compte opérateur) \\
\rowcolor{trowa}
POST & \texttt{/auth/login} & Public & Connexion, renvoie le jeton JWT \\
GET, PATCH & \texttt{/auth/me} & Connecté & Lire et modifier son profil \\
\rowcolor{trowa}
POST & \texttt{/auth/me/password} & Connecté & Changer son mot de passe \\
GET, POST & \texttt{/auth/users} & Admin & Lister et créer des comptes \\
\rowcolor{trowa}
PATCH, DELETE & \texttt{/auth/users/\{id\}} & Admin & Modifier et supprimer un compte \\
PATCH & \texttt{/auth/users/\{id\}/toggle-active} & Admin & Activer ou désactiver un compte \\
\rowcolor{trowa}
POST & \texttt{/auth/users/\{id\}/reset-password} & Admin & Réinitialiser un mot de passe \\
GET, DELETE & \texttt{/auth/audit-logs} & Admin & Consulter et purger le journal d'audit \\
\multicolumn{4}{@{}l}{\textit{Module transmission (détection)}} \\
\rowcolor{trowa}
GET & \texttt{/transmission/kpis} & Public & Chiffres clés de la page d'accueil \\
GET & \texttt{/transmission/report} & Op., Admin & Rapport d'alerte par type d'attaque \\
\rowcolor{trowa}
GET & \texttt{/transmission/history} & Op., Admin & Historique des entraînements \\
DELETE & \texttt{/transmission/history/\{id\}} & Admin & Supprimer un entraînement \\
\rowcolor{trowa}
GET, POST, DELETE & \texttt{/transmission/dataset} & Admin & Lire, importer (CSV) et retirer le jeu importé \\
POST & \texttt{/transmission/train} & Admin & Lancer l'entraînement en arrière-plan \\
\rowcolor{trowa}
GET & \texttt{/transmission/training-status} & Op., Admin & Messages de progression depuis le dernier lu \\
GET & \texttt{/transmission/training-target} & Admin & Modèle retenu et son réglage d'alerte \\
\multicolumn{4}{@{}l}{\textit{Module etude (lecture des résultats et temps réel)}} \\
\rowcolor{trowa}
GET & \texttt{/overview} & Op., Admin & Indicateurs de la vue d'ensemble \\
GET & \texttt{/etude/metrics}, \texttt{/etude/runs} & Op., Admin & Métriques des modèles, liste des runs \\
\rowcolor{trowa}
GET & \texttt{/etude/timeline} & Op., Admin & Alertes et attaques dans le temps \\
GET & \texttt{/etude/shap} & Op., Admin & Importance SHAP globale et par scénario \\
\rowcolor{trowa}
GET & \texttt{/etude/dataset} & Op., Admin & Indicateurs du jeu de données \\
GET & \texttt{/etude/live/info} & Op., Admin & Modèle, seuil et vitesse par défaut du rejeu \\
\rowcolor{trowa}
POST & \texttt{/etude/live/start} & Op., Admin & Démarrer une lecture côté serveur \\
GET & \texttt{/etude/live/\{id\}/next} & Op., Admin & Paquet suivant de mesures (1 à 200) \\
\rowcolor{trowa}
DELETE & \texttt{/etude/live/\{id\}} & Op., Admin & Arrêter la lecture \\
\multicolumn{4}{@{}l}{\textit{Module parametres}} \\
GET & \texttt{/parametres} & Op., Admin & Lire les paramètres, bornes et valeurs par défaut \\
\rowcolor{trowa}
PUT & \texttt{/parametres} & Admin & Modifier les paramètres (tracé dans l'audit) \\
\hline
\end{tabularx}
\end{table}

\subsection{Données persistantes}

Tout ce que le tableau de bord affiche est lu en base ; aucun résultat n'est
écrit en dur. Les tables se regroupent en trois familles :
\begin{itemize}
  \item \textbf{Gestion} : \texttt{users}, \texttt{audit\_logs},
        \texttt{app\_settings} ;
  \item \textbf{Mesures et événements} : \texttt{transmission\_dataset}
        (mesures et variables), \texttt{transmission\_episodes} (épisodes
        d'attaque), \texttt{transmission\_predictions},
        \texttt{transmission\_alerts} ;
  \item \textbf{Résultats d'entraînement} : \texttt{transmission\_runs} (un
        enregistrement par entraînement), \texttt{transmission\_model\_results},
        \texttt{transmission\_metrics}, \texttt{transmission\_shap},
        \texttt{transmission\_learning\_curve}.
\end{itemize}

\subsection{Rejeu temps réel et entraînement}

Deux traitements s'exécutent côté serveur et sont suivis par interrogation
périodique du client.

\begin{itemize}
  \item \textbf{Rejeu temps réel.} \texttt{POST /etude/live/start} crée une
        lecture des mesures de la période choisie (entraînement, validation ou
        test). Le tableau de bord réclame ensuite chaque seconde, avec
        \texttt{GET /etude/live/\{id\}/next}, un paquet de mesures dont la
        taille correspond à la vitesse. Le serveur calcule le score du modèle,
        applique le seuil et la règle de persistance, puis renvoie les points,
        les débuts et fins d'alerte et les compteurs (alertes vraies et
        fausses).
  \item \textbf{Entraînement.} \texttt{POST /transmission/train} lance
        l'entraînement en arrière-plan ; un second lancement simultané est
        refusé (code 409). Le modèle retenu par l'étude de référence est
        réentraîné avec son réglage d'alerte, sur le jeu importé ou sur celui du
        projet. Le client lit les messages de progression avec
        \texttt{GET /transmission/training-status}, puis le résultat est
        enregistré en base et devient consultable dans « Résultats ».
\end{itemize}

% =====================================================================
\section{Réalisation du tableau de bord}

Le tableau de bord est une application monopage React~19 écrite en TypeScript
et construite avec Vite. Nous ne décrivons que les fichiers effectivement
chargés par l'application (tableau~\ref{tab:front}).

\begin{table}[H]
\centering
\small
\renewcommand{\arraystretch}{1.4}
\setlength{\arrayrulewidth}{0.8pt}
\arrayrulecolor{thead}
\caption{Fichiers utilisés du tableau de bord (\texttt{frontend/src}).}
\label{tab:front}
\begin{tabularx}{\textwidth}{@{}>{\raggedright\arraybackslash}p{4.6cm} >{\raggedright\arraybackslash}X@{}}
\rowcolor{thead}
\textcolor{white}{\textbf{Fichier}} & \textcolor{white}{\textbf{Rôle}} \\
\hline
\texttt{main.tsx}, \texttt{App.tsx} &
Point d'entrée et routage : \texttt{HashRouter}, fournisseurs de thème, de
session et d'infobulles, routes publiques (\texttt{/}, \texttt{/login},
\texttt{/register}) et route protégée donnant accès au tableau de bord. \\
\rowcolor{trowa}
\texttt{components/PrivateRoute.tsx} &
Redirige vers la connexion si l'utilisateur n'est pas authentifié. \\
\texttt{layouts/DashboardLayout.tsx} &
Coquille du tableau de bord : menu en trois groupes filtré selon le rôle,
barre supérieure (thème, menu du compte), garde d'accès par page (« Accès
refusé ») et chargement différé des pages. \\
\rowcolor{trowa}
\texttt{contexts/AuthContext.tsx}, \texttt{lib/auth.ts} &
Session : connexion, déconnexion, jeton JWT conservé et vérifié à l'échéance,
en-têtes d'authentification des requêtes. \\
\texttt{contexts/ThemeContext.tsx} &
Thème clair ou sombre. \\
\rowcolor{trowa}
\texttt{components/etude/kit.tsx} &
Boîte à outils commune des pages : accès à l'API avec cache (\texttt{useApi}),
en-têtes de page, sections, indicateurs, tableaux, palette des graphiques et
libellés français des modèles et des scénarios. \\
\texttt{components/AuthShell.tsx}, \texttt{BandMark.tsx} &
Gabarit des pages de connexion et d'inscription ; logo. \\
\rowcolor{trowa}
\texttt{components/ui/} (tooltip, toast, toaster, use-toast, dropdown-menu) &
Composants Radix UI utilisés : infobulles, notifications et menu du compte. \\
\texttt{lib/features.ts}, \texttt{lib/utils.ts}, \texttt{types.ts} &
Regroupement des variables du modèle par capteur, type et fenêtre ;
utilitaire de classes CSS ; type du thème. \\
\hline
\end{tabularx}
\end{table}

Le tableau~\ref{tab:pages} donne, pour chaque page, l'accès par rôle et les
routes de l'API qu'elle appelle. Les pages de la zone privée sont chargées à
la demande (\textit{lazy loading}), ce qui allège le premier affichage. Le
temps réel et le suivi de l'entraînement fonctionnent par interrogation
périodique (une requête par seconde) et non par connexion persistante.

\begin{table}[H]
\centering
\small
\renewcommand{\arraystretch}{1.4}
\setlength{\arrayrulewidth}{0.8pt}
\arrayrulecolor{thead}
\caption{Pages du tableau de bord, accès et routes de l'API appelées.}
\label{tab:pages}
\begin{tabularx}{\textwidth}{@{}>{\raggedright\arraybackslash}p{3.6cm} >{\centering\arraybackslash}p{1.9cm} >{\raggedright\arraybackslash}X@{}}
\rowcolor{thead}
\textcolor{white}{\textbf{Page}} & \textcolor{white}{\textbf{Accès}} &
\textcolor{white}{\textbf{Routes de l'API}} \\
\hline
\multicolumn{3}{@{}l}{\textit{Pages publiques}} \\
\texttt{HomePage} & Tous & \texttt{GET /transmission/kpis} \\
\rowcolor{trowa}
\texttt{LoginPage}, \texttt{RegisterPage} & Tous &
\texttt{POST /auth/login}, \texttt{POST /auth/register} \\
\multicolumn{3}{@{}l}{\textit{Surveillance}} \\
\rowcolor{trowa}
\texttt{OverviewPage} & Op., Admin & \texttt{GET /overview} \\
\texttt{TempsReelPage} & Op., Admin &
\texttt{GET /etude/live/info}, \texttt{POST /etude/live/start},
\texttt{GET /etude/live/\{id\}/next}, \texttt{DELETE /etude/live/\{id\}} \\
\rowcolor{trowa}
\texttt{RapportPage} & Op., Admin & \texttt{GET /transmission/report} \\
\multicolumn{3}{@{}l}{\textit{Analyse du modèle}} \\
\texttt{ResultatsPage} & Op., Admin & \texttt{GET /etude/metrics}, \texttt{GET /etude/runs} \\
\rowcolor{trowa}
\texttt{DetectionPage} & Op., Admin & \texttt{GET /etude/timeline} \\
\texttt{ShapPage} & Op., Admin & \texttt{GET /etude/shap} \\
\rowcolor{trowa}
\texttt{DatasetPage} & Op., Admin & \texttt{GET /etude/dataset} \\
\multicolumn{3}{@{}l}{\textit{Administration}} \\
\texttt{EntrainementPage} & Admin &
\texttt{POST /transmission/train}, \texttt{GET /transmission/training-status},
\texttt{GET /transmission/training-target}, \texttt{GET},
\texttt{POST}, \texttt{DELETE /transmission/dataset},
\texttt{GET}, \texttt{DELETE /transmission/history} \\
\rowcolor{trowa}
\texttt{ParametresPage} & Admin & \texttt{GET}, \texttt{PUT /parametres} \\
\texttt{AdminUsersPage} & Admin & \texttt{/auth/users} (liste, création, modification,
activation, réinitialisation du mot de passe, suppression) \\
\rowcolor{trowa}
\texttt{AuditLogsPage} & Admin & \texttt{GET}, \texttt{DELETE /auth/audit-logs} \\
\multicolumn{3}{@{}l}{\textit{Compte}} \\
\texttt{ProfilePage} & Op., Admin & \texttt{PATCH /auth/me}, \texttt{POST /auth/me/password} \\
\hline
\end{tabularx}
\end{table}

\subsection{Aperçu de l'interface}

Les figures suivantes présentent les pages principales. Placez les captures
d'écran dans le dossier \texttt{figures/} avec les noms indiqués ; tant
qu'elles sont absentes, un cadre les remplace.

\newcommand{\capture}[3]{%
  \begin{figure}[H]\centering
  \IfFileExists{figures/#1.png}%
    {\includegraphics[width=0.92\textwidth]{figures/#1.png}}%
    {\fbox{\parbox[c][5cm][c]{0.8\textwidth}{\centering Capture à insérer : \texttt{figures/#1.png}}}}
  \caption{#2}\label{#3}
  \end{figure}}

\capture{vue_ensemble}{Page « Vue d'ensemble » : dernier entraînement et détection par scénario.}{fig:ui-overview}
\capture{temps_reel}{Page « Temps réel » : mesures, score, alertes et attaques réelles côte à côte.}{fig:ui-live}
\capture{detection}{Page « Détection vs réalité » : alertes du modèle et épisodes d'attaque dans le temps.}{fig:ui-detection}
\capture{rapport}{Page « Rapport d'alerte » : consigne à appliquer par type d'attaque.}{fig:ui-report}
\capture{shap}{Page « Explicabilité SHAP » : importance globale et par scénario.}{fig:ui-shap}
\capture{entrainement}{Page « Entraînement » (administrateur) : import du CSV et suivi en direct.}{fig:ui-train}

% =====================================================================
\section{Vérification par les tests}

Le backend est vérifié par 34 tests automatisés pytest, répartis en six
fichiers (tableau~\ref{tab:tests}). Ils s'exécutent dans le conteneur avec
\texttt{docker compose exec app python -m pytest tests -q}.

\begin{table}[H]
\centering
\small
\renewcommand{\arraystretch}{1.4}
\setlength{\arrayrulewidth}{0.8pt}
\arrayrulecolor{thead}
\caption{Tests automatisés du backend.}
\label{tab:tests}
\begin{tabularx}{\textwidth}{@{}>{\raggedright\arraybackslash}p{4.8cm} >{\centering\arraybackslash}p{1.4cm} >{\raggedright\arraybackslash}X@{}}
\rowcolor{thead}
\textcolor{white}{\textbf{Fichier}} & \textcolor{white}{\textbf{Tests}} &
\textcolor{white}{\textbf{Ce qui est vérifié}} \\
\hline
\texttt{test\_transmission\_protocol} & 13 &
Protocole d'évaluation : découpage 70/15/15 par blocs entiers, 20 rotations,
zone de garde, aucun épisode ni fenêtre de variables à cheval sur deux rôles,
variables causales, règle d'alerte $k$ sur $n$, choix du modèle sur la
validation seule. \\
\rowcolor{trowa}
\texttt{test\_transmission} & 12 &
Rapport pour l'opérateur : consigne par famille d'attaque, niveau d'alerte,
chiffres publics de la page d'accueil, entraînement web avec le modèle de
référence. \\
\texttt{test\_etude\_roles} & 3 &
Étude de référence stockée et lue en base ; refus (403) d'un opérateur sur les
comptes, les paramètres et les entraînements ; rejeu temps réel (seuil,
persistance, alertes vraies et fausses). \\
\rowcolor{trowa}
\texttt{test\_transmission\_db} & 3 &
Un entraînement terminé est stocké avec une ligne par modèle ; les échecs
gardent leur statut ; historique et résultats lus en base. \\
\texttt{test\_overview} & 2 &
Base vide : blocs vides ; base remplie : chiffres conformes à la base. \\
\rowcolor{trowa}
\texttt{test\_settings} & 1 &
Origines CORS lues depuis la configuration. \\
\hline
\end{tabularx}
\end{table}

% =====================================================================
\section{Déploiement}

L'application se déploie en une commande, \texttt{docker compose up -d --build}
(figure~\ref{fig:deploiement}). L'image est construite en deux étapes : une
étape Node compile le tableau de bord avec Vite, puis une image Python embarque
l'API, les modèles et le tableau de bord compilé. L'API sert donc l'interface et
les routes REST sur un seul port (8000), accessible en local uniquement et à
publier derrière un proxy inverse. Le processus s'exécute sous un utilisateur
non privilégié, avec un contrôle de santé. Au premier démarrage, le jeu de
données est reconstruit à partir des traces brutes dans un volume.

\begin{figure}[H]
\centering
\begin{tikzpicture}[
  box/.style={draw=thead, fill=tlayer, rounded corners=3pt, align=center,
              font=\footnotesize, minimum height=1.2cm, text width=3.8cm, inner sep=4pt},
  vol/.style={draw=gray!60, fill=gray!10, rounded corners=3pt, align=center,
              font=\scriptsize, minimum height=0.8cm, text width=2.4cm},
  arr/.style={-Latex, thick, draw=thead},
  lbl/.style={font=\scriptsize\itshape, text=gray!80!black}
]
\node[box] (nav) at (0,0) {\textbf{Navigateur}\\Tableau de bord};
\node[box] (app) at (5.2,0) {\textbf{Service \texttt{app}}\\FastAPI + tableau de bord\\port 8000};
\node[box] (db)  at (10.4,0) {\textbf{Service \texttt{db}}\\PostgreSQL 16};
\node[vol] (v1)  at (4.2,-2.1) {volume \texttt{data}};
\node[vol] (v2)  at (6.4,-2.1) {volume \texttt{models}};
\node[vol] (v3)  at (10.4,-2.1) {volume \texttt{pgdata}};
\node[box, fill=gray!10, draw=gray!60] (ad) at (10.4,2.1) {\textbf{Service \texttt{adminer}}\\Consultation de la base\\port 8080};
\draw[arr] (nav) -- node[lbl, above] {HTTP} (app);
\draw[arr] (app) -- node[lbl, above] {SQL asynchrone} (db);
\draw[arr] (app) -- (v1);
\draw[arr] (app) -- (v2);
\draw[arr] (db) -- (v3);
\draw[arr, gray] (ad) -- (db);
\end{tikzpicture}
\caption{Déploiement de VigiloChain avec Docker Compose.}
\label{fig:deploiement}
\end{figure}

% =====================================================================
\section{Traçabilité des besoins}

Le tableau~\ref{tab:trace} relie chaque besoin fonctionnel à la page du
tableau de bord et aux routes de l'API qui le réalisent.

\begin{table}[H]
\centering
\footnotesize
\renewcommand{\arraystretch}{1.35}
\setlength{\arrayrulewidth}{0.8pt}
\arrayrulecolor{thead}
\caption{Matrice de traçabilité des besoins fonctionnels.}
\label{tab:trace}
\begin{tabularx}{\textwidth}{@{}>{\bfseries}p{1.1cm} >{\raggedright\arraybackslash}p{3.4cm} >{\raggedright\arraybackslash}X@{}}
\rowcolor{thead}
\textcolor{white}{\textbf{Besoin}} & \textcolor{white}{\textbf{Page}} &
\textcolor{white}{\textbf{Routes de l'API}} \\
\hline
BF1 & \texttt{OverviewPage} & \texttt{GET /overview} \\
\rowcolor{trowa}
BF2 & \texttt{TempsReelPage} & \texttt{/etude/live/info}, \texttt{start}, \texttt{\{id\}/next}, \texttt{DELETE \{id\}} \\
BF3 & \texttt{DetectionPage} & \texttt{GET /etude/timeline} \\
\rowcolor{trowa}
BF4 & \texttt{RapportPage} & \texttt{GET /transmission/report} \\
BF5 & \texttt{ResultatsPage} & \texttt{GET /etude/metrics}, \texttt{GET /etude/runs} \\
\rowcolor{trowa}
BF6 & \texttt{ShapPage} & \texttt{GET /etude/shap} \\
BF7 & \texttt{DatasetPage} & \texttt{GET /etude/dataset} \\
\rowcolor{trowa}
BF8 & \texttt{EntrainementPage} & \texttt{/transmission/dataset}, \texttt{train}, \texttt{training-status}, \texttt{training-target}, \texttt{history} \\
BF9 & \texttt{ParametresPage} & \texttt{GET}, \texttt{PUT /parametres} \\
\rowcolor{trowa}
BF10 & \texttt{AdminUsersPage} & \texttt{/auth/users} et sous-routes \\
BF11 & \texttt{AuditLogsPage} & \texttt{GET}, \texttt{DELETE /auth/audit-logs} \\
\rowcolor{trowa}
BF12 & \texttt{HomePage}, \texttt{LoginPage}, \texttt{RegisterPage} & \texttt{/transmission/kpis}, \texttt{/auth/login}, \texttt{/auth/register} \\
BF13 & \texttt{ProfilePage}, \texttt{DashboardLayout} & \texttt{PATCH /auth/me}, \texttt{POST /auth/me/password} \\
\hline
\end{tabularx}
\end{table}

% =====================================================================
\section*{Conclusion}
\addcontentsline{toc}{section}{Conclusion}
Ce chapitre a présenté la réalisation de VigiloChain : un backend FastAPI qui
protège chaque route par rôle et lit ses résultats en base, un tableau de bord
React dont chaque page correspond à un besoin fonctionnel, 34 tests
automatisés et un déploiement en une commande avec Docker Compose. La matrice
de traçabilité montre que les treize besoins du chapitre~\ref{chap2} sont tous
couverts. Le chapitre suivant conclut le rapport et ouvre des perspectives.
